\documentclass[12pt,a4paper,openany]{memoir}

\usepackage[utf8]{inputenc}
\usepackage[spanish,es-nodecimaldot]{babel}
%\usepackage{showframe}
\usepackage{lipsum}
\usepackage{microtype}
\usepackage{kantlipsum}
\usepackage{graphicx}
\usepackage{hyperref}

% La portada (?)
% La fuente
\usepackage[adobe-utopia]{mathdesign}
\usepackage[T1]{fontenc}

% Márgenes y dónde va el bloque de texto
% \medievalpage
% \isopage
% \semiisopage
\setlrmarginsandblock{1in}{1in}{*}
\setulmarginsandblock{2in}{2.5in}{*}
\checkandfixthelayout

% Encabezado y pie de página
%\pagestyle{headings}
\nouppercaseheads
\makepagestyle{nuestro}
\makeevenfoot{nuestro}{\thepage}{}{}
\makeoddfoot{nuestro}{}{}{\thepage}
\makeevenhead{nuestro}{\theauthor}{}{\leftmark}
\makeoddhead{nuestro}{\rightmark}{}{\today}
\makeheadrule{nuestro}{\textwidth}{1pt}
%\makefootrule{nuestro}{0.5\textwidth}{1pt}{10pt}
\pagestyle{nuestro}

% El estilo del encabezado de un capítulo
%\chapterstyle{wilsondob}
% section
% hangnum
% companion
% article
% reparticle
% demo
% bianchi
% bringhurst
% brotherton
% chappell
% culver
% dash
% ell
% lyhne
% veelo
% verville
% crosshead
% komalike
%\chapterstyle{culver}
\makechapterstyle{AlexanderGrebenkov}{%
    \renewcommand\chapterheadstart{\vspace*{\beforechapskip}\hrule\medskip}\renewcommand\chapnamefont{\normalfont\large\scshape}
    \renewcommand\chapnumfont{\normalfont\large\scshape}
    \renewcommand\chaptitlefont{\normalfont\large\scshape}
    \renewcommand\printchaptername{\hfill\S}
    \renewcommand\chapternamenum{ }
    \renewcommand\printchapternum{\chapnumfont \thechapter}
    \renewcommand\afterchapternum{. }
    \renewcommand\printchapternonum{\hfill}
    \renewcommand\afterchaptertitle{
        \par\nobreak\medskip\hrule\vskip\afterchapskip
    }
}
\chapterstyle{AlexanderGrebenkov}

\makeatletter
\newcommand\thickhrulefill{\leavevmode \leaders \hrule height 1ex \hfill \kern \z@}
\setlength\midchapskip{10pt}
\makechapterstyle{VZ14}{
    \setlength{\beforechapskip}{2\baselineskip}
    \renewcommand\chapterheadstart{\vspace*{\beforechapskip}}
    \renewcommand\chapternamenum{\space}
    \renewcommand\printchaptername{\space}
    \renewcommand\chapnamefont{\Large\scshape}
    \renewcommand\printchapternum{%
        \chapnamefont\null\thickhrulefill\quad
        \@chapapp\space\thechapter\quad\thickhrulefill
    }
    \renewcommand\printchapternonum{%
        \par\thickhrulefill\par\vskip\midchapskip
        \hrule\vskip\midchapskip
    }
    \renewcommand\chaptitlefont{\Huge\scshape\centering}\renewcommand\afterchapternum{%
        \par\nobreak\vskip\midchapskip
        \hrule\vskip\midchapskip
    }
    \renewcommand\afterchaptertitle{%
        \par\vskip\midchapskip
        \hrule\nobreak\vskip\afterchapskip
    }
}
\makeatother
%\chapterstyle{VZ14}
%\renewcommand{\clearforchapter}{}
\setbeforesecskip{3.5ex plus 1ex minus .2ex}
\setsecheadstyle{\normalfont\normalsize\scshape}
\setaftersecskip{2.3ex plus .2ex}
\setsubsecheadstyle{\normalfont\normalsize\scshape}

\author{inm}
\title{El título del trabajo}
\date{\today}

\begin{document}

% Una opción para hacer una portada es usar titlingpage
%\begin{titlingpage}
%    \maketitle
%\end{titlingpage}

\thispagestyle{empty}

\setlength{\unitlength}{1in}
\begin{picture}(0,0)
    \put(0,0) {\includegraphics[width=1in]{figs/ciencias.png}}
\end{picture}
\setlength{\unitlength}{1pt}

\begin{minipage}[c][0.8\textheight][c]{1in}
    \hspace{0.45in}
    \rule{1pt}{0.8\textheight}
\end{minipage}
\hspace{0.1\textwidth}
\begin{minipage}[c][0.8\textheight][c]{0.5\textwidth}
    \centering
    {\HUGE\bfseries \thetitle}\\[3\baselineskip]
    {\Large\itshape Un bello subtítulo}\\[2\baselineskip]
    {\large inm \\[\baselineskip] anónimo \\[\baselineskip] sujeto 1 \\[\baselineskip] holi }
    \vfill
\end{minipage}

\vfill

\setlength{\unitlength}{1in}
\begin{picture}(0,0)
    \put(0,0) {\includegraphics[width=1in]{figs/unam.png}}
\end{picture}
\setlength{\unitlength}{1pt}


\clearpage

{
    \renewcommand\cleardoublepage\newpage
    \frontmatter
}

\tableofcontents*

{
    \renewcommand\cleardoublepage\newpage
    \mainmatter
}

\chapter{Hola}

\noindent \lipsum[1-6]

\section{Ja, ja}

\noindent Chuchu\footnote{Soy un pie de página}\footnote{A ver} \lipsum[7-12]

\chapter{Mundo}

\noindent \kant[1-4]

\section{Lo que sea}

\noindent \kant[5-7]

\end{document}